
\chapter{Sistemas Populacionais}
\label{cap:10}

A célula isolada pode ser modelada em um único tubo de ensaio. A espécie
distribuída em um lago, em uma floresta ou em uma metrópole, não. Assim que se
transfere a lente da bioquímica para a ecologia, a entidade a ser estudada deixa
de ser uma molécula e passa a ser uma \emph{população} --- um agregado de
indivíduos que nascem, morrem, migram e interagem em ambientes heterogêneos. A
descrição matemática dessa dinâmica é uma das áreas mais antigas e mais
férteis da biologia teórica, com raízes que vão de Malthus (1798) ao trabalho
contemporâneo em metapopulações e epidemiologia matemática.

A pergunta central deste capítulo é simples de enunciar e difícil de responder:
como variam as densidades de populações ao longo do tempo e do espaço, e como
essa variação pode ser prevista, controlada ou restaurada? Para respondê-la,
discutiremos cinco tipos de modelos, em ordem crescente de complexidade. Os
modelos de \textbf{crescimento de uma única espécie} introduzirão as ferramentas
fundamentais --- equações diferenciais ordinárias, equações de diferença,
análise de estabilidade local. Os modelos de \textbf{interação entre espécies}
mostrarão como predadores, presas, competidores e mutualistas se acoplam em
sistemas dinâmicos. Os \textbf{modelos epidemiológicos} aplicarão as mesmas
técnicas a doenças infecciosas, onde o indivíduo é a "espécie" e as
transições são eventos (S $\to$ I $\to$ R) em vez de nascimentos e mortes. Por
fim, a \textbf{dinâmica espacial} incorporará difusão, metapopulações e padrões
de Turing --- fechando o capítulo com a observação de que a estrutura espacial
é tão importante quanto a interação local.

\section{Fundamentos da Dinâmica Populacional}
\label{sec:10-fundamentos}

\begin{definicaoenv}[Sistema Populacional]
Sistema populacional é o conjunto de indivíduos de uma mesma espécie (ou de
espécies interagentes) cuja densidade varia no tempo e no espaço em resposta a
processos elementares de natalidade, mortalidade, migração e interação. A
descrição matemática dessa variação é o objeto da ecologia teórica e da
epidemiologia matemática.
\end{definicaoenv}

A escrita matemática mais simples para uma população $N(t)$ de uma única
espécie é o balanço entre quatro processos demográficos:
\begin{equation}
\frac{dN}{dt} = \underbrace{b N}_{\text{nascimento}}
               - \underbrace{d N}_{\text{morte}}
               + \underbrace{I}_{\text{imigração}}
               - \underbrace{E}_{\text{emigração}}
\end{equation}
onde $b$ e $d$ são taxas per capita e $I$, $E$ são taxas absolutas. Em
contextos onde nascimento e morte são os únicos processos relevantes, $I = E$
e a equação se reduz à forma malthusiana. Quando os recursos do ambiente são
limitados, o termo de natalidade é modificado por uma função de saturação, e a
equação adquire a forma logística. Estudaremos ambos os regimes a seguir.

\subsection{Crescimento Exponencial: O Modelo de Malthus}
\label{sub:image10-exponencial}

O modelo mais simples possível de dinâmica populacional postula que a taxa de
crescimento é proporcional ao tamanho da população:
\begin{equation}
\frac{dN}{dt} = r N
\end{equation}

O parâmetro $r = b - d$ é a \emph{taxa intrínseca de crescimento per capita} --- a
diferença entre as taxas de natalidade e mortalidade, ambas expressas na
mesma escala por indivíduo. A solução é imediata:
\begin{equation}
N(t) = N_0\, e^{rt}
\end{equation}

O modelo exponencial descreve com precisão surpreendente populações em regimes
de baixa densidade, longe da capacidade de suporte do ambiente. Bactérias em
meio fresco, populações de insetos durante surtos e processos tumorais
estabelecidos obedecem a essa cinética durante a fase exponencial. O
\emph{tempo de duplicação}, $t_d = \ln 2 / r$, é uma grandeza útil para
diagnóstico: para a \textit{Escherichia coli} em condições ideais, $t_d \approx
20$ min; para uma cultura de células humanas, $t_d \approx 24$ h; para a
população humana global nas décadas de 1960--1970, $r \approx 0{,}02\;\mathrm{a}^{-1}$
e $t_d \approx 35$ anos.

O modelo exponencial também descreve \emph{decrescimento} quando $r < 0$.
Populações em ambientes hostis --- como populações pequenas confinadas a
\emph{habitats} degradados --- seguem decaimento exponencial. O destino do
\emph{Panthera tigris altaica}, por exemplo, depende crucialmente de que $r$
atinja valores positivos, o que requer ações integradas de proteção do
\emph{habitat}, antipoaching e manejo genético.

\subsection{Crescimento Logístico: Introduzindo a Capacidade de Suporte}
\label{sub:image10-logistico}

O modelo exponencial ignora que nenhum recurso é ilimitado. Em 1838, o
matemático belga Pierre-François Verhulst propôs uma modificação simples ---
introduzir um termo de inibição dependente da densidade, que reduz o
crescimento à medida que a população se aproxima da \emph{capacidade de suporte}
$K$ (a densidade máxima sustentável pelo ambiente):
\begin{equation}
\frac{dN}{dt} = r N\left(1 - \frac{N}{K}\right)
\end{equation}

A solução analítica é a clássica sigmoidal:
\begin{equation}
N(t) = \frac{K}{1 + \left(\dfrac{K - N_0}{N_0}\right)\, e^{-rt}}
\end{equation}

A densidade $N = K/2$ merece destaque: é o ponto de inflexão da sigmoidal e o
ponto em que a taxa $dN/dt$ é máxima. Populações abaixo desse valor estão em
fase de aceleração; acima dele, em fase de desaceleração. Modelos logísticos
descrevem bem populações de protozoários em cultura fechada, crescimento
tumoral em fase precoce (modelo de Gompertz, primo do logístico) e, com
adaptações, populações de peixes em reservatórios.

A implementação computacional do logístico é imediata usando integradores
padrão. O trecho de código a seguir reproduz a curva sigmoidal clássica:

\begin{lstlisting}[language=Python, caption=Modelo logistico]
import numpy as np
import matplotlib.pyplot as plt
from scipy.integrate import odeint

def logistic(N, t, r, K):
    return r * N * (1 - N/K)

r, K, N0 = 0.8, 1000, 50
t = np.linspace(0, 20, 500)
N = odeint(logistic, N0, t, args=(r, K))

plt.figure(figsize=(8, 5))
plt.plot(t, N, 'b-', lw=2, label='N(t)')
plt.axhline(y=K, color='r', ls='--', lw=2, label=f'K = {K}')
plt.xlabel('Tempo'); plt.ylabel('Populacao')
plt.title('Crescimento Logistico')
plt.legend(); plt.grid(alpha=0.3)
plt.tight_layout()
plt.show()
\end{lstlisting}

\subsection{Efeito Allee: A Maldição das Populações Pequenas}
\label{sub:image10-allee}

Em muitas espécies --- especialmente vertebrados sociais --- densidades
pequenas reduzem, em vez de aumentar, a taxa de crescimento per capita. As
causas vão da dificuldade de encontrar parceiros sexuais à redução da defesa
coletiva contra predadores. Warder Clyde Allee documentou este fenômeno na
década de 1930, e a equação que o carrega é:
\begin{equation}
\frac{dN}{dt} = r N \left(1 - \frac{N}{K}\right) \left(\frac{N}{A} - 1\right)
\end{equation}

O parâmetro $A$, com $0 < A < K$, é a \emph{densidade crítica} de Allee --- a
menor densidade compatível com crescimento populacional positivo. A equação
admite três pontos de equilíbrio: $N^* = 0$ (extinção, estável), $N^* = A$
(instável) e $N^* = K$ (estável). A interpretação é direta: uma população
iniciada entre $0$ e $A$ colapsa irreversivelmente; apenas populações
iniciadas acima de $A$ sobrevivem. O gráfico do crescimento em função de $N$
cruza o eixo em três pontos, formando a "curva em J invertido" que é assinatura
do efeito Allee forte.

A implicação para a conservação é profunda. Populações pequenas estão
inerentemente expostas a \emph{extinção por depopulação}: flutuações
estocásticas podem levá-las abaixo do limiar $A$, de onde não há retorno.
Programas bem-sucedidos de recuperação --- como o do \textit{Arabidopsis
arenosa} na Suíça ou da águia-pescadora no Reino Unido --- exigem população
inicial grande o suficiente para garantir densidade acima do limiar de Allee.

\subsection{Modelos em Tempo Discreto: Ricker e Beverton-Holt}
\label{sub:image10-discreto}

As equações diferenciais descrevem populações com gerações que se sobrepõem.
Quando as gerações são discretas --- insetos anuais, peixes em habitats
sazonais, populações microbianas em ciclos sincronizados --- a equação natural é
uma \emph{equação de diferença}. O modelo de Ricker é emblemático:
\begin{equation}
N_{t+1} = N_t\, e^{r\left(1 - \frac{N_t}{K}\right)}
\end{equation}

Apesar de sua forma inocente, o modelo de Ricker produz comportamento
surpreendentemente rico: convergência monotônica para $K$ em baixos valores de
$r$; oscilações amortecidas em valores intermediários; oscilações sustentadas;
ciclos de período 2, 4, 8 ...; e caos para $r$ suficientemente alto (a partir
de $r \approx 2{,}69$ para $K = 100$).

O modelo alternativo é o de Beverton-Holt:
\begin{equation}
N_{t+1} = \frac{R\, N_t}{1 + (R - 1)\, N_t / K}
\end{equation}
onde $R$ é a taxa de reprodução líquida. Diferentemente do modelo de Ricker,
o Beverton-Holt sempre converge monotonicamente a $K$ --- não exibe oscilações
nem caos.

A análise gráfica de sistemas discretos é feita no plano de fases
$(N_t, N_{t+1})$ --- o chamado \emph{diagrama cobweb}. A população evolui
seguindo o gráfico da função $f(N_t)$ e refletindo-se a cada iteração na
linha de identidade $N_{t+1} = N_t$. Os pontos fixos são as interseções
entre $f$ e a linha de identidade; sua estabilidade é dada pela derivada $|f'(N^*)| < 1$.

O trecho de código a seguir ilustra o efeito dramático de $r$ sobre a
dinâmica do modelo de Ricker:

\begin{lstlisting}[language=Python, caption=Dinamica de Ricker para varios r]
import numpy as np
import matplotlib.pyplot as plt

def ricker(N, r, K):
    return N * np.exp(r * (1 - N/K))

K = 100
N0 = 10
t_max = 50
fig, axes = plt.subplots(1, 3, figsize=(15, 4))

for i, r in enumerate([1.5, 2.0, 2.7]):
    N = np.zeros(t_max); N[0] = N0
    for t in range(t_max - 1):
        N[t+1] = ricker(N[t], r, K)
    axes[i].plot(N, 'b.-', ms=4)
    axes[i].axhline(y=K, color='r', ls='--', label=f'K = {K}')
    axes[i].set_xlabel('Geracao t'); axes[i].set_ylabel('N(t)')
    axes[i].set_title(f'r = {r}')
    axes[i].legend(); axes[i].grid(alpha=0.3)

plt.tight_layout()
plt.show()
\end{lstlisting}


\section{Interações entre Espécies}
\label{sec:10-interacoes}

A descrição de uma população isolada é apenas o ponto de partida. Na natureza,
organismos de diferentes espécies interagem --- competem por alimento e espaço,
predam ou são predados, cooperam ou parasitam. A classificação clássica
dessas interações usa a notação de sinais: $(+)$ benéfico, $(-)$ prejudicial,
$(0)$ neutro. Competição é $(-,-)$: ambas as espécies perdem. Predação e
parasitismo são $(+,-)$. Mutualismo é $(+,+)$. Comensalismo é $(+,0)$ ---
uma espécie se beneficia enquanto a outra é indiferente. Os modelos
matemáticos destas interações generalizam os de uma única espécie, com a
complexidade adicional de que múltiplos equilíbrios podem existir e sua
estabilidade depende dos parâmetros relativos das espécies envolvidas.

\subsection{Lotka-Volterra: A Dinâmica Predador-Presa}
\label{sub:image10-lotka-volterra}

O modelo mais antigo e influente de interação entre espécies foi formulado
independentemente por Alfred Lotka (1925) e Vito Volterra (1926). Para um
sistema com presas $V(t)$ e predadores $P(t)$ em interação, a hipótese central
é que encontros seguem a lei de ação das massas (proporcionais a $VP$), e que
a presença de presas beneficia o crescimento dos predadores proporcionalmente
à taxa de encontro:
\begin{align}
\frac{dV}{dt} &= \alpha V - \beta V P \\
\frac{dP}{dt} &= \delta \beta V P - \gamma P
\end{align}

O parâmetro $\alpha$ é a taxa intrínseca de crescimento das presas na ausência
de predação; $\beta$ é a taxa de encontro/predação; $\delta$ é a eficiência
com que presas consumidas convertem-se em novos predadores; e $\gamma$ é a
taxa de mortalidade dos predadores.

Análise. O sistema tem dois equilíbrios triviais --- $(0,0)$ e o que resulta
em extinção de uma espécie --- e um equilíbrio não-trivial de coexistência:
\begin{equation}
V^* = \frac{\gamma}{\delta \beta}, \qquad P^* = \frac{\alpha}{\beta}
\end{equation}

O Jacobiano no equilíbrio tem traço zero e determinante positivo. Seus
autovalores são puramente imaginários $\pm i \sqrt{\alpha \gamma}$, indicando
que o equilíbrio é um \emph{centro} --- nem atrator nem repulsor. Trajetórias
próximas formam órbitas fechadas ao redor do equilíbrio. A interpretação
biológica é célebre: \emph{o sistema exibe oscilações perpétuas} cuja amplitude
e período dependem das condições iniciais, mas não se amortecem. Limitando o
modelo (introduzindo capacidade de suporte para as presas, predação
dependente de densidade ou termo de inibição do predador) obtém-se ciclos
limites estáveis --- que é o que se observa em sistemas reais como o clássico
caso do \emph{Lynx canadensis} e da lebre-americana (\textit{Lepus americanus})
registrado pelos dados de peles da Hudson Bay Company.

\begin{lstlisting}[language=Python, caption=Sistema predador-presa de Lotka-Volterra]
import numpy as np
import matplotlib.pyplot as plt
from scipy.integrate import odeint

def lotka_volterra(y, t, alpha, beta, delta, gamma):
    V, P = y
    return [alpha*V - beta*V*P,
            delta*beta*V*P - gamma*P]

alpha, beta, delta, gamma = 1.0, 0.1, 0.75, 1.5
y0 = [40, 9]
t = np.linspace(0, 50, 1000)
sol = odeint(lotka_volterra, y0, t,
             args=(alpha, beta, delta, gamma))

V_eq, P_eq = gamma/(delta*beta), alpha/beta

fig, axes = plt.subplots(1, 2, figsize=(13, 5))
axes[0].plot(t, sol[:, 0], 'b-', lw=2, label='Presas (V)')
axes[0].plot(t, sol[:, 1], 'r-', lw=2, label='Predadores (P)')
axes[0].set_xlabel('Tempo'); axes[0].set_ylabel('Densidade')
axes[0].legend(); axes[0].grid(alpha=0.3)

axes[1].plot(sol[:, 0], sol[:, 1], 'g-', lw=2)
axes[1].plot(V_eq, P_eq, 'r*', ms=15, label='Equilibrio')
axes[1].axhline(P_eq, ls='--', color='r', alpha=0.5, label='dV/dt=0')
axes[1].axvline(V_eq, ls='--', color='b', alpha=0.5, label='dP/dt=0')
axes[1].set_xlabel('Presas (V)'); axes[1].set_ylabel('Predadores (P)')
axes[1].legend(); axes[1].grid(alpha=0.3)
plt.tight_layout(); plt.show()
\end{lstlisting}

\subsection{Competição e o Princípio de Exclusão}
\label{sub:image10-competicao}

Lotka e Volterra, no mesmo estudo, propuseram o modelo de competição interespecífica:
\begin{align}
\frac{dN_1}{dt} &= r_1 N_1 \left(1 - \frac{N_1 + \alpha_{12} N_2}{K_1}\right) \\
\frac{dN_2}{dt} &= r_2 N_2 \left(1 - \frac{N_2 + \alpha_{21} N_1}{K_2}\right)
\end{align}

O parâmetro $\alpha_{12}$ expressa o efeito de um indivíduo da espécie 2
sobre a espécie 1, em unidades de "equivalentes" da espécie 1; $\alpha_{21}$
é a grandeza simétrica. Os quatro regimes qualitativos possíveis são definidos
por combinações das desigualdades entre $\alpha_{12} K_2$, $K_1$, $\alpha_{21}
K_1$ e $K_2$:

\begin{itemize}
\item Se $K_1 > \alpha_{12} K_2$ e $K_2 > \alpha_{21} K_1$: coexistência
  estável em equilíbrio interior.
\item Se $K_1 > \alpha_{12} K_2$ e $K_2 < \alpha_{21} K_1$: espécie 1 vence.
\item Se $K_1 < \alpha_{12} K_2$ e $K_2 > \alpha_{21} K_1$: espécie 2 vence.
\item Se $K_1 < \alpha_{12} K_2$ e $K_2 < \alpha_{21} K_1$: o vencedor
  depende das condições iniciais --- bistabilidade.
\end{itemize}

Este último caso é a base experimental do \emph{princípio da exclusão
competitiva}, atribuído a Gause: duas espécies em competição por um único
recurso limitante não podem coexistir indefinidamente. O famoso experimento
de Gause (1934) com \textit{Paramecium aurelia} e \textit{Paramecium caudatum}
em cultura com a mesma fonte alimentar mostrou exatamente o previsto --- uma
espécie exterminava a outra. Hutchinson reinterpretou o resultado em termos de
\emph{nicho ecológico}: a coexistência prolongada exige diferenciação de nicho
(recurso, espaço, tempo).

\subsection{Mutualismo e Simbiose}
\label{sub:image10-mutualismo}

Quando a interação é $(+,+)$, ambos os parceiros se beneficiam. O modelo
canônico generaliza o logístico substituindo a capacidade de suporte fixa
por uma função crescente da densidade do parceiro:
\begin{align}
\frac{dN_1}{dt} &= r_1 N_1 \left(1 - \frac{N_1}{K_1 + \alpha_{12} N_2}\right) \\
\frac{dN_2}{dt} &= r_2 N_2 \left(1 - \frac{N_2}{K_2 + \alpha_{21} N_1}\right)
\end{align}

Os mutualismos \emph{facultativos} --- aqueles em que cada espécie pode
sobreviver sozinha --- diferem dos mutualismos \emph{obrigatórios}, nos quais
a sobrevivência individual depende de interações com a outra espécie. As
micorrizas, que trocam carbono da planta por fósforo e nitrogênio do solo, são
um caso quase obrigatório. A polinização por abelhas é tipicamente
facultativa para as plantas (vento, aves e mariposas também polinizam), mas
obrigatória para muitas culturas agrícolas --- fato que justifica programas
de proteção de polinizadores em níveis mundiais.

\subsection{Modelo de Nicholson-Bailey: Hospedeiro-Parasitóide}
\label{sub:image10-nicholson-bailey}

A interação entre hospedeiros e parasitóides (insetos cujas larvas se
desenvolvem dentro de hospedeiros, eventualmente matando-os) é uma das
aplicações mais elegantes de dinâmica discreta. O modelo de Nicholson-Bailey
assume distribuição de Poisson para os ataques:
\begin{align}
H_{t+1} &= \lambda H_t e^{-a P_t} \\
P_{t+1} &= c H_t \left(1 - e^{-a P_t}\right)
\end{align}

O termo exponencial $e^{-a P_t}$ é a fração de hospedeiros que escapa do
parasitismo na geração $t$; $\lambda$ é a taxa de reprodução do hospedeiro;
$a$ é a área de descoberta do parasitóide; $c$ é o número médio de parasitóides
emergindo por hospedeiro parasitado. O modelo prevê que, na ausência de
respostas funcionais ou numéricas, a interação é sempre instável --- o
parasitóide superpredar seu hospedeiro, levando ambos à extinção. Mecanismos
de regulação (dependência de densidade, refugos espaciais, alternância de
gerações) são necessários para explicar a coexistência estável observada.

Esse modelo é central no \emph{controle biológico de pragas}: a introdução
controlada de parasitóides específicos (vespas do gênero \textit{Trichogramma}
contra lagartas, por exemplo) é uma alternativa ambientalmente amigável a
inseticidas químicos --- mas depende crucialmente do ajuste fino dos parâmetros
do modelo à biologia da praga-alvo.


\section{Modelos Epidemiológicos}
\label{sec:10-epidemiologia}

A epidemiologia matemática é, em certo sentido, uma aplicação particular de
dinâmica populacional: em vez de espécies biológicas, as "espécies" são
\emph{estados de saúde}, e os parâmetros de nascimento e morte são substituídos
por transições entre compartimentos. Essa mudança de perspectiva, formalizada
por Kermack e McKendrick em 1927, transformou a saúde pública: hoje,
projeções de surtos e avaliação de estratégias de vacinação dependem
inteiramente de modelos dessa classe.

\subsection{Comportamentalização por Compartimentos}
\label{sub:image10-sir}

A abordagem padrão particiona a população em classes discretas segundo seu
status em relação à doença:
\begin{itemize}
\item \textbf{S} (\emph{Susceptíveis}): indivíduos que podem ser infectados.
\item \textbf{E} (\emph{Expostos}): infectados em período de incubação, ainda
  não infecciosos.
\item \textbf{I} (\emph{Infectados}): infectados e capazes de transmitir.
\item \textbf{R} (\emph{Recuperados}): imunes por recuperação natural ou
  vacinação.
\item \textbf{D} (\emph{Mortos}): na versão que inclui mortalidade pela doença.
\end{itemize}

A passagem de um compartimento a outro é regida por taxas que dependem do
contágio ($\beta$, mediado pelo contato infectante) e da recuperação ($\gamma$,
inversa do período infeccioso). A transmissão segue a lei de ação das massas:
o número de contatos infectantes por unidade de tempo é proporcional ao
produto $S \cdot I$ e à densidade de encontros.

\subsection{O Modelo SIR Clássico}
\label{sub:image10-sir-classico}

O modelo mais simples de uma doença conferindo imunidade permanente é o
\emph{SIR}:
\begin{align}
\frac{dS}{dt} &= -\beta \frac{SI}{N} \\
\frac{dI}{dt} &= \beta \frac{SI}{N} - \gamma I \\
\frac{dR}{dt} &= \gamma I
\end{align}

onde $N = S + I + R$ é mantido constante (caso de doença com mortalidade
desprezível em escala epidêmica). O parâmetro $\beta$ é a taxa de transmissão
(contatos infectantes por unidade de tempo), e $\gamma$ é a taxa de
recuperação. A escala natural de tempo é o período infeccioso médio, dado
por $1/\gamma$.

\subsection{Número Básico de Reprodução $R_0$}
\label{sub:image10-R0}

\begin{definicaoenv}[Número básico de reprodução $R_0$]
$R_0$ é o número médio de infecções secundárias causadas por um único
indivíduo infectado introduzido em uma população inteiramente suscetível.
É o parâmetro mais importante da epidemiologia matemática: governa o limiar
entre extinção e propagação da epidemia, e sintetiza o efeito combinado de
transmissibilidade, frequência de contatos e duração da infecciosidade.
\end{definicaoenv}

Para o modelo SIR simples, $R_0$ é o produto de duas grandezas operacionais:
\begin{equation}
R_0 = \frac{\beta}{\gamma} = \underbrace{\beta}_{\text{contatos infectantes
por unidade de tempo}} \times \underbrace{\frac{1}{\gamma}}_{\text{duração do
período infeccioso}}
\end{equation}

O limiar epidêmico é exatamente $R_0 = 1$:
\begin{itemize}
\item $R_0 < 1$: a epidemia se extingue --- cada infectado gera em média
  menos de um caso secundário.
\item $R_0 = 1$: endemia estável --- o sistema se mantém em equilíbrio
  exatamente limítrofe.
\item $R_0 > 1$: epidemia se propaga.
\end{itemize}

Valores típicos de $R_0$ para doenças humanas ilustram a diversidade do
problema: sarampo $R_0 \approx 12$--$18$, varíola $R_0 \approx 5$--$7$,
rubéola $R_0 \approx 5$--$7$, pólio $R_0 \approx 5$--$7$, COVID-19 (variante
original) $R_0 \approx 2$--$3$, influenza sazonal $R_0 \approx 1$--$2$,
ebola $R_0 \approx 1{,}5$--$2{,}5$. Uma doença mais transmissível requer
controle mais agressivo.

\subsection{Threshold de Vacinação}
\label{sub:image10-vaccination}

A vacinação move indivíduos de $S$ para $R$ sem passar por $I$. Se uma fração
$f$ da população é vacinada, a densidade efetiva de suscetíveis reduz-se pelo
fator $(1 - f)$, e o parâmetro de controle efetivo torna-se
$R_0^{\mathrm{eff}} = R_0(1 - f)$. Para evitar epidemia precisamos
$R_0^{\mathrm{eff}} < 1$, donde:
\begin{equation}
f > 1 - \frac{1}{R_0}
\end{equation}

Para sarampo ($R_0 \approx 15$), $f > 0{,}93$ --- 93\% de cobertura. Esta é a
exigência que torna a hesitação vacinal tão perigosa: a queda de 5\% na
cobertura vacinal em algumas comunidades brasileiras nas últimas décadas já é
suficiente para permitir surtos --- e os surtos aparecem.

\begin{exemploenv}[Threshold do sarampo]
O sarampo requer cobertura acima de 93\% para imunidade de rebanho. A
queda global de cobertura para cerca de 80\% entre 2020 e 2022 --- durante a
pandemia de COVID-19 --- criou o cenário ideal para o ressurgimento que
ocorreu em 2023--2024, com surtos em vários continentes, incluindo o Brasil.
\end{exemploenv}

\subsection{Modelo SIS: Doenças sem Imunidade Duradoura}
\label{sub:image10-sis}

Nem todas as doenças conferem imunidade permanente. Resfriados comuns,
infecções sexualmente transmissíveis como gonorreia e clamídia, e várias
infecções intestinais reaparecem nos mesmos indivíduos após recuperação.
Nesses casos, a população retorna de $I$ diretamente a $S$, em vez de a $R$:
\begin{align}
\frac{dS}{dt} &= -\beta \frac{SI}{N} + \gamma I \\
\frac{dI}{dt} &= \beta \frac{SI}{N} - \gamma I
\end{align}

Quando $R_0 > 1$, o sistema atinge equilíbrio endêmico com $I^* = N(1 - 1/R_0)$
e $S^* = N/R_0$. A doença persiste indefinidamente na população --- é o regime
\emph{endêmico}, distinto do regime epidêmico transitório do SIR.

\subsection{Modelo SEIR: Com Período de Incubação}
\label{sub:image10-seir}

Para doenças com período de incubação significativo --- Ebola, COVID-19,
infecções por HIV, tuberculose --- o modelo SIR subestima a latência entre
infecção e infecciosidade. O modelo \emph{SEIR} adiciona um compartimento
explícito $E$ para "expostos não infecciosos":
\begin{align}
\frac{dS}{dt} &= -\beta \frac{SI}{N} \\
\frac{dE}{dt} &= \beta \frac{SI}{N} - \sigma E \\
\frac{dI}{dt} &= \sigma E - \gamma I \\
\frac{dR}{dt} &= \gamma I
\end{align}

O parâmetro $\sigma$ é a taxa de progressão de $E$ para $I$ --- seu inverso
é o período de incubação médio. Embora $R_0$ permaneça formalmente
$R_0 = \beta/\gamma$, a \emph{dinâmica temporal} difere: a fase de propagação
inicial é mais lenta (a doença tem que "esperar" a incubação), mas a curva
de pico é similar. Em COVID-19, o modelo SEIR modificado com demografia,
estrutura etária e intervenções não-farmacêuticas (distanciamento, máscaras,
fechamento de escolas) foi a base das projeções que guiaram políticas
sanitárias em 2020--2022.

A implementação computacional do SIR ilustra como extrair medidas
epidemiológicas básicas --- pico de infectados, tempo do pico, taxa de ataque
final, limiar epidêmico --- da solução numérica do sistema:

\begin{lstlisting}[language=Python, caption=Modelo SIR e sua analise]
import numpy as np
import matplotlib.pyplot as plt
from scipy.integrate import odeint

def sir_model(y, t, beta, gamma, N):
    S, I, R = y
    return [-beta*S*I/N,
             beta*S*I/N - gamma*I,
             gamma*I]

N, beta, gamma = 10000, 0.5, 0.1
R0 = beta/gamma
I0, S0, R0_ini = 10, N - 10, 0
t = np.linspace(0, 200, 1000)
sol = odeint(sir_model, [S0, I0, R0_ini], t,
             args=(beta, gamma, N))
S, I, R = sol.T

fig, axes = plt.subplots(1, 2, figsize=(13, 5))
axes[0].plot(t, S, 'b-', lw=2, label='Suscetiveis (S)')
axes[0].plot(t, I, 'r-', lw=2, label='Infectados (I)')
axes[0].plot(t, R, 'g-', lw=2, label='Recuperados (R)')
axes[0].set_xlabel('Tempo (dias)'); axes[0].set_ylabel('Individuos')
axes[0].set_title(f'Dinamica SIR ($R_0$ = {R0:.1f})')
axes[0].legend(); axes[0].grid(alpha=0.3)

peak_t = t[np.argmax(I)]
axes[1].plot(S, I, 'purple', lw=2)
axes[1].axvline(x=N/R0, color='red', ls='--',
                label=f'S = N/R0 = {N/R0:.0f}')
axes[1].set_xlabel('Suscetiveis (S)'); axes[1].set_ylabel('Infectados (I)')
axes[1].set_title('Retrato de Fase'); axes[1].legend(); axes[1].grid(alpha=0.3)
print(f"Pico: {I.max():.0f} individuos em t = {peak_t:.1f} dias")
print(f"Taxa de ataque final: {R[-1]/N*100:.1f}%")
plt.tight_layout(); plt.show()
\end{lstlisting}


\section{Dinâmica Espacial de Populações}
\label{sec:10-espaco}

Os modelos até aqui descreveram populações espacialmente homogêneas --- a
densidade $N(t)$ é a mesma em todos os pontos do ambiente. Na realidade,
\emph{não existe o espaço médio}. Diferenças locais em disponibilidade de
alimento, abrigo e refúgio contra predadores produzem gradientes espaciais
que influenciam profundamente a dinâmica populacional. Esta seção introduz
os modelos canônicos de dinâmica espacial: metapopulações e equações de
reação-difusão, com destaque para ondas viajantes e padrões de Turing.

\subsection{Metapopulações: O Modelo de Levins}
\label{sub:image10-metapop}

\begin{definicaoenv}[Metapopulação]
Metapopulação é um conjunto de populações locais (chamadas \emph{patches})
conectadas por dispersão, no qual cada população local pode sofrer extinção
e ser recolonizada por imigrantes de outros patches. O conceito foi
formalizado por Levins em 1969 a partir da observação de que habitats
naturais são intrinsecamente fragmentados.
\end{definicaoenv}

O modelo de Levins trata uma \emph{fração} $p(t)$ dos patches como ocupada
e a fração $(1 - p)$ como vazia. As transições são:
\begin{itemize}
\item \textbf{Colonização}: patches vazios são ocupados na taxa $c \cdot p(1 - p)$,
  proporcional à fração ocupada (fonte de propágulos) e à fração vazia
  (alvos disponíveis).
\item \textbf{Extinção local}: patches ocupados são perdidos na taxa $e \cdot p$.
\end{itemize}

A equação resultante é simples e elegante:
\begin{equation}
\frac{dp}{dt} = c p (1 - p) - e p = p\,\bigl(c(1 - p) - e\bigr)
\end{equation}

O equilíbrio não-trivial $p^* = 1 - e/c$ é positivo (e portanto a
metapopulação persiste) se e somente se $c > e$. A condição \emph{c $>$ e
$}$ é uma das formulações mais elegantes em ecologia: a colonização deve
superar a extinção local para que a metapopulação persista. Em sua ausência,
a fração de habitats ocupados colapsa a zero.

O modelo básico de Levins pode ser enriquecido de várias maneiras: patches
heterogêneos (tamanho variável, qualidade diferente), \emph{corredores}
ecológicos (conectividade entre fragmentos), efeito resgate (\emph{rescue effect},
em que imigrantes reduzem a probabilidade de extinção local), e dinâmica
intra-patch mais detalhada (cada patch seguindo seu próprio modelo de
crescimento, com migração entre patches). A implementação com redes
heterogêneas de patches é discutida na seção final.

\subsection{Equações de Reação-Difusão: Fisher-KPP}
\label{sub:image10-fisher-kpp}

A próxima classe de modelos pressupõe um espaço contínuo e utiliza a equação
de reação-difusão:
\begin{equation}
\frac{\partial N}{\partial t} = \underbrace{f(N)}_{\text{reação local}}
                              + \underbrace{D \nabla^2 N}_{\text{dispersão}}
\end{equation}

onde $f(N)$ é a dinâmica local de crescimento (logística, por exemplo, ou uma
forma qualquer) e $D$ é o coeficiente de difusão, que descreve o movimento
browniano --- a versão mais simples de dispersão. Em 1D, o Laplaciano é a
segunda derivada $\partial^2 N/\partial x^2$.

A equação emblemática para invasões biológicas é a equação de
\emph{Fisher-KPP}, que combina crescimento logístico com difusão:
\begin{equation}
\frac{\partial N}{\partial t} = r N \left(1 - \frac{N}{K}\right)
                              + D \frac{\partial^2 N}{\partial x^2}
\end{equation}

Fisher (1937) e Kolmogorov-Petrovskii-Piskunov (1937) mostraram
independentemente que essa equação admite soluções de \emph{onda viajante} ---
perfis de densidade $N(x - ct)$ que se propagam com velocidade constante
mantendo a forma. A velocidade mínima (selecionada pelas condições iniciais
suaves) é:
\begin{equation}
c_{\min} = 2\sqrt{r D}
\end{equation}

A equação Fisher-KPP é um dos pilares da teoria de invasão. Foi aplicada
com sucesso à propagação da lebre europeia (\textit{Oryctolagus
cuniculagus}) na Austrália, à expansão de QTL de resistência a inseticidas
em populações de \textit{Anopheles}, e à dinâmica de ondas em frentes de
queimadas em ecossistemas secos.

A simulação numérica da equação em 1D usa diferenças finitas centrais para o
Laplaciano. Um trecho de código típico segue o formato abaixo:

\begin{lstlisting}[language=Python, caption=Equacao de Fisher-KPP em 1D]
import numpy as np
import matplotlib.pyplot as plt

def fisher_kpp_1d(N, r, K, D, dx):
    reaction  = r * N * (1 - N/K)
    diffusion = np.zeros_like(N)
    diffusion[1:-1] = D * (N[2:] - 2*N[1:-1] + N[:-2]) / dx**2
    diffusion[0] = diffusion[-1] = 0  # Neumann: sem fluxo
    return reaction + diffusion

L = 50; nx = 200
x = np.linspace(0, L, nx); dx = x[1] - x[0]
r, K, D = 0.5, 1.0, 0.1

N = np.zeros(nx)
N[nx//4:nx//4+10] = K  # pulso inicial

t_max, dt = 100, 0.01
steps = int(t_max / dt)
for _ in range(steps):
    N += dt * fisher_kpp_1d(N, r, K, D, dx)

plt.figure(figsize=(9, 4))
plt.plot(x, N, 'b-', lw=2)
plt.axhline(K, color='r', ls='--', label='K')
plt.xlabel('x'); plt.ylabel('N'); plt.legend(); plt.grid(alpha=0.3)
plt.tight_layout(); plt.show()
\end{lstlisting}

\subsection{Padrões de Turing: A Origem Química da Forma}
\label{sub:image10-turing}

Em 1952, Alan Turing publicou um dos artigos mais influentes da biologia
teórica --- \emph{The Chemical Basis of Morphogenesis}. Nele, Turing mostrou
que um sistema de duas espécies reagentes com difusão diferenciada pode
desestabilizar espontaneamente o estado homogêneo, gerando padrões espaciais
estacionários. Os ingredientes são dois: dois morfógenos --- um ativador (com
curto alcance) e um inibidor (com longo alcance) --- e difusão mais rápida do
inibidor do que do ativador:
\begin{align}
\frac{\partial u}{\partial t} &= f(u, v) + D_u \nabla^2 u \\
\frac{\partial v}{\partial t} &= g(u, v) + D_v \nabla^2 v
\end{align}

A análise de estabilidade linear ao redor do equilíbrio homogêneo mostra que,
embora o equilíbrio seja estável na ausência de difusão ($f$ e $g$ são
estabilizantes), ele se torna instável quando $D_v \gg D_u$. A partir do
estado homogêneo, o sistema evolui espontaneamente para padrões de manchas
(\emph{spots}) ou listras. O comprimento de onda característico desses
padrões é:
\begin{equation}
\lambda \sim 2\pi \sqrt{\frac{D_v}{\mu}}
\end{equation}
onde $\mu$ é a taxa de crescimento linearizada do modo mais instável.

Biólogos foram lentos para abraçar a ideia, mas hoje há forte evidência de
que padrões de Turing estão subjacentes à organização de estômatos em
folhas, aos folículos de pena em aves e à pelagem de certos peixes e
felinos. A matemática, aplicada a sistemas biológicos concretos, ajudou a
transformar uma especulação teórica em um programa de pesquisa vivo.

\subsection{Metapopulações em Rede}
\label{sub:image10-metapop-network}

A formalização mais geral trata a metapopulação como uma rede de patches em
que cada um segue sua própria dinâmica local e as trocas migratórias entre
patches vizinhos são mediadas por coeficientes de acoplamento. Para um patch
$i$, a equação generalizada é:
\begin{equation}
\frac{d N_i}{dt} = f(N_i) + \sum_{j=1}^{n} m_{ij}\,(N_j - N_i)
\end{equation}
onde $m_{ij}$ é a taxa de migração entre o patch $j$ e o patch $i$, e
$f(N_i)$ é a dinâmica local (por exemplo, logística).

A topologia da rede determina o sincronismo da metapopulação: alta
conectividade ($\sum_j m_{ij}$ grande) tende a sincronizar oscilações entre
patches, \emph{aumentando} o risco de extinção global em sistemas com
estabilização dependente de densidade; baixa conectividade gera dinâmicas
independentes, com \emph{efeito de portfolio} análogo ao da diversificação em
finanças. As topologias que têm recebido mais atenção teórica são grade
regular, rede aleatória Erdős-Rényi, redes \emph{scale-free} e grafos do
tipo mundo-pequeno --- cada uma com propriedades próprias de sincronização
e vulnerabilidade.

\begin{exemploenv}[Corredores ecológicos]
A aplicação prática mais direta de metapopulações em rede é a conservação
biológica em paisagens fragmentadas. Corredores ecológicos conectando
fragmentos de habitat --- matas ciliares restauradas, sistemas de árvores
pontilhando lavouras, passagens de fauna sobre rodovias --- são projetados
para aumentar a conectividade efetiva $m_{ij}$ da metapopulação de uma
espécie-alvo. Aumentos modestos de $m_{ij}$ em paisagens críticas para a
mantença de predadores de topo (onças-pintadas, lobos-guará) podem ser
decisivos para a sobrevivência populacional a longo prazo.
\end{exemploenv}


\section{Exercícios}
\label{sec:10-exercicios}

\begin{exercicio}
\textbf{Crescimento logístico analítico.} Considere uma população seguindo o
modelo logístico com $r = 0{,}6\;\mathrm{dia}^{-1}$ e $K = 5000$.
Calcule (a) $N(10)$ quando $N(0) = 500$; (b) o tempo necessário para a
população atingir 90\% de $K$; (c) a taxa máxima $dN/dt$ e a densidade em que
ela ocorre.
\end{exercicio}

\begin{exercicio}
\textbf{Efeito Allee.} O modelo de Allee
$\tfrac{dN}{dt} = rN(1 - N/K)(N/A - 1)$
é considerado com $r = 0{,}5$, $K = 1000$, $A = 100$.
(a) Identifique os pontos de equilíbrio e analise sua estabilidade local.
(b) Implemente o modelo em Python e simule para $N(0) = 50$ e $N(0) = 150$.
(c) Explique as implicações para programas de conservação que trabalham
com populações pequenas.
\end{exercicio}

\begin{exercicio}
\textbf{Bifurcação no mapa de Ricker.} Implemente o modelo de Ricker e
explore a transição de comportamento à medida que $r$ varia de 0 a 3.
(a) Construa um diagrama de bifurcação plotando os valores de $N$ visitados
em $t = 200$ (após descartar o transiente) como função de $r$.
(b) Identifique o início de oscilações de período 2 e de comportamento
caótico.
(c) Compare com o mapa de Beverton-Holt no mesmo intervalo de $r$ ---
quantitativamente, o que difere?
\end{exercicio}

\begin{exercicio}
\textbf{Análise do Lotka-Volterra.} Para o sistema predador-presa com
$\alpha = 1{,}0$, $\beta = 0{,}1$, $\delta = 0{,}75$, $\gamma = 1{,}5$:
(a) Encontre o equilíbrio não-trivial;
(b) Calcule a Jacobiana no equilíbrio e seus autovalores;
(c) Com base nos autovalores, argumente sobre o tipo de estabilidade;
(d) Implemente em Python e plote séries temporais de $V(t)$ e $P(t)$ e o
retrato de fase.
\end{exercicio}

\begin{exercicio}
\textbf{Competição interespecífica.} Implemente o sistema de competição de
Lotka-Volterra em Python. (a) Simule os quatro cenários
(uma espécie vence, outra vence, coexistência, exclusão dependente de CI)
usando combinações adequadas de $K_1$, $K_2$, $\alpha_{12}$, $\alpha_{21}$.
(b) Construa um retrato de fase com isoclinas para cada caso.
(c) Adicione ruído gaussiano às derivadas e investigue como flutuações
ambientais modificam o destino competitivo.
\end{exercicio}

\begin{exercicio}
\textbf{Modelo SIR e gestão de saúde.} Uma doença tem $\beta = 0{,}4\;\mathrm{dia}^{-1}$
e $\gamma = 0{,}1\;\mathrm{dia}^{-1}$ em uma cidade de $N = 10^4$ habitantes.
Inicialmente $I(0) = 20$. (a) Calcule $R_0$. (b) Simule 200 dias e
identifique o pico de infectados e a data em que ocorre. (c) Estime a taxa
de ataque final. (d) Determine a fração mínima de vacinação para prevenir a
epidemia. (e) Repita os itens (b)-(d) para variantes com $\beta = 0{,}6$ e
$\beta = 0{,}8$ e discuta o que muda em termos de política sanitária.
\end{exercicio}

\begin{exercicio}
\textbf{Modelo SEIR e período de incubação.} Modifique seu código SIR para
incluir um compartimento $E$ com $\sigma = 0{,}2\;\mathrm{dia}^{-1}$
(período médio de incubação de 5 dias). Mantenha $\beta = 0{,}5$ e $\gamma = 0{,}1$.
(a) Compare o tempo do pico de infectados no SIR e no SEIR.
(b) Verifique se o $R_0$ muda; explique por que sim ou por que não.
(c) Investigue como o atraso entre exposição e infecciosidade altera a
\emph{fase} da curva epidêmica.
\end{exercicio}

\begin{exercicio}
\textbf{Metapopulação de Levins.} O modelo de Levins
$\tfrac{dp}{dt} = cp(1-p) - ep$
descreve a dinâmica da fração de patches ocupados em uma paisagem. (a)
Considere uma paisagem onde a colonização local ocorre a $c = 0{,}5\;\mathrm{ano}^{-1}$
e a extinção local a $e = 0{,}1\;\mathrm{ano}^{-1}$. Qual a fração de equilíbrio
ocupado? (b) Como esse equilíbrio se altera se uma rodovia interrompe $30\%$
dos corredores ecológicos (reduzindo $c$ em $30\%$)? (c) Quaisquer que sejam
os valores, a velocidade de retorno ao equilíbrio é proporcional a $(c-e)$.
Discuta por que a velocidade de resposta da metapopulação a uma mudança de
\emph{habitat} depende crucialmente da margem $c - e$.
\end{exercicio}

\begin{exercicio}
\textbf{Modelo de Nicholson-Bailey.} Implemente o sistema parasitóide-hospedeiro
em Python. (a) Mostre que, com os parâmetros padrão
$\lambda = 2$, $a = 0{,}02$, $c = 1$, a interação leva à extinção de ambas
as espécies em poucas gerações, em concordância com a análise de estabilidade
do modelo original. (b) Adicione uma resposta funcional do tipo II ao parasitóide
(substitua o termo de encontro $a P_t$ por $a P_t / (1 + a h P_t)$ com
$h = 0{,}5$). Discuta como esse ingrediente estabilizador pode levar à
coexistência. (c) Encontre, por simulação, um par de parâmetros
$(\lambda, a, c, h)$ que produz ciclos persistentes de hospedeiro e parasitóide.
\end{exercicio}

\begin{exercicio}
\textbf{Simulação de Fisher-KPP.} Implemente a equação de Fisher-KPP em 1D
usando o esquema de diferenças finitas centrais. (a) Plote o perfil de $N(x)$
em três instantes ($t = 0$, $50$, $100$) e meça a velocidade da onda
viajando por ajuste linear do ponto de meia-altura. Compare com o valor
analítico $c_{\min} = 2\sqrt{r D}$. (b) Repita a simulação para diferentes
valores de $D$ e verifique a relação $c_{\min} \propto \sqrt{D}$.
(c) Varie a condição inicial (em vez de um pulso compacto, use uma rampa
suave) e discuta por que a velocidade converge para $c_{\min}$ quando a
condição inicial é monótona.
\end{exercicio}

\begin{exercicio}
\textbf{Projeto integrador: dinâmica espacial de uma invasão biológica.}
Considere um caso real --- fictício ou documentado --- de introdução de uma
espécie invasora em uma paisagem fragmentada. (a) Formule um modelo
matemático em 2D que combine crescimento logístico, difusão e efeitos de
\emph{habitat} heterogêneo (representado por uma função $K(x,y)$ espacialmente
variável). (b) Implemente em Python usando diferenças finitas em uma grade
$100 \times 100$. (c) Plote mapas de densidade em função do tempo e identifique
\emph{corredores} de invasão. (d) Proponha alterações no uso da terra que
minimizem a taxa de avanço. (e) Compare seus resultados com observações
empíricas de uma invasão real (\textit{Lantana camara} em cerrado, \textit{Brachypodium sylvaticum} em florestas de
Washington, etc.).
\end{exercicio}

\section{Notas Bibliográficas}
\label{sec:10-bibliografia}

A ecologia teórica clássica está consolidada em três livros texto de uso
corrente. \textit{Mathematical Biology I: An Introduction} de James D. Murray
(3ª edição, Springer) é o trabalho de referência unificado em inglês --- cobre
desde modelos de uma única espécie até ondas viajantes, autômatos celulares e
o problema da propagação do HIV. \textit{Mathematical Models in Biology} de Leah
Edelstein-Keshet (SIAM) é mais acessível e bi-Centrado, com exemplos extensos
de osciladores gênicos, propagação de ondas em meios excitáveis e dinâmica
de populações. \textit{A Biologist's Guide to Mathematical Modeling in Ecology
and Evolution} de Sarah Otto e Toby Day (Princeton) é indicado para leitores
que preferem começar com questões biológicas e chegar gradualmente à matemática;
inclui discussão profunda de modelos evolutivos, teoria de jogos e inferência
em ecologia.

A \emph{epidemiologia matemática} tem dois textos seminais que todo profissional
de saúde pública deveria consultar. \textit{Infectious Diseases of Humans:
Dynamics and Control} de Roy Anderson e Robert May (Oxford) é o livro de
referência para modelagem de doenças infecciosas humanas --- formaliza o SIR
e suas extensões, discute estruturação etária, redes de contato e a
epidemiologia do HIV/AIDS. \textit{Modeling Infectious Diseases in Humans
and Animals} de Matt Keeling e Pejman Rohani (Princeton) é mais moderno e
computacionalmente orientado, com código em MATLAB e extensões para doenças
zoonóticas. \textit{Mathematical Models in Population Biology and
Epidemiology} de Fred Brauer e Carlos Castillo-Chávez é a apresentação mais
sistemática disponível, ideal para quem quiser dominar o tratamento
matemático completo.

Para os tópicos avançados --- metapopulações, difusão, padrões de Turing ---
são úteis \textit{An Illustrated Guide to Theoretical Ecology} de Ted Case, com
ênfase em intuição e exemplos; \textit{Elements of Mathematical Ecology} de
Mark Kot; e \textit{Resource Competition and Community Structure} de David
Tilman, leitura indispensável para quem estuda ecologia de comunidades. Os
artigos seminais --- Malthus (1798), Verhulst (1838), Lotka (1925),
Volterra (1926), Nicholson e Bailey (1935), Levins (1969), Turing (1952),
Fisher (1937), Kolmogorov et al.\ (1937), Kermack e McKendrick (1927) ---
continuam recomendados para a formação do historiador da disciplina.

Recursos computacionais incluem o pacote estatístico \texttt{EpiModel} (em R),
amplamente usado em epidemiologia, e a ferramenta \texttt{PyDSTool} (Python)
para análise de sistemas dinâmicos. \texttt{COPASI} é uma plataforma
multiplataforma GUI para simulação e análise de modelos de sistemas
biológicos, especialmente útil em cursos introdutórios.

\begin{thebibliography}{99}
\bibitem{murray} Murray J.D. (2002). \textit{Mathematical Biology I: An
Introduction}, 3ª ed. Springer, Nova York.

\bibitem{edelstein} Edelstein-Keshet L. (2005). \textit{Mathematical Models
in Biology}. SIAM, Philadelphia.

\bibitem{otto} Otto S.P., Day T. (2007). \textit{A Biologist's Guide to
Mathematical Modeling in Ecology and Evolution}. Princeton University Press,
Princeton.

\bibitem{keeling} Keeling M.J., Rohani P. (2008). \textit{Modeling Infectious
Diseases in Humans and Animals}. Princeton University Press, Princeton.

\bibitem{anderson} Anderson R.M., May R.M. (1991). \textit{Infectious Diseases
of Humans: Dynamics and Control}. Oxford University Press, Oxford.

\bibitem{brauer} Brauer F., Castillo-Chavez C. (2012). \textit{Mathematical
Models in Population Biology and Epidemiology}, 2ª ed. Springer, Nova York.

\bibitem{kermack} Kermack W.O., McKendrick A.G. (1927). A contribution to the
mathematical theory of epidemics. \textit{Proceedings of the Royal Society
of London A}, 115:700--721.

\bibitem{lotka} Lotka A.J. (1925). \textit{Elements of Physical Biology}.
Williams \& Wilkins, Baltimore.

\bibitem{volterra} Volterra V. (1926). Fluctuations in the abundance of a
species considered mathematically. \textit{Nature}, 118:558--560.

\bibitem{nicholson} Nicholson A.J., Bailey V.A. (1935). The balance of animal
populations. \textit{Proceedings of the Zoological Society of London},
1:551--598.

\bibitem{levins} Levins R. (1969). Some demographic and genetic consequences
of environmental heterogeneity for biological control. \textit{Bulletin of
the Entomological Society of America}, 15:237--240.

\bibitem{turing} Turing A.M. (1952). The chemical basis of morphogenesis.
\textit{Philosophical Transactions of the Royal Society B}, 237:37--72.

\bibitem{fisher} Fisher R.A. (1937). The wave of advance of advantageous
genes. \textit{Annals of Eugenics}, 7:355--369.

\bibitem{tilman} Tilman D. (1982). \textit{Resource Competition and Community
Structure}. Princeton University Press, Princeton.

\bibitem{kot} Kot M. (2001). \textit{Elements of Mathematical Ecology}.
Cambridge University Press, Cambridge.

\bibitem{case} Case T.J. (2000). \textit{An Illustrated Guide to Theoretical
Ecology}. Oxford University Press, Oxford.

\bibitem{epimodel} Jenness S., Goodreau S.M., Morris M. (2018). EpiModel: an
R package for mathematical modeling of infectious disease over networks.
\textit{Journal of Statistical Software}, 84:1--47. Disponível em
\url{https://epimodel.org}.

\bibitem{copasi} Hoops S., Sahle S., Gauges R., Lee C., Pahle J., Simus N.,
Singhal M., Xu L., Mendes P., Kummer U. (2006). COPASI --- a COmplex PAthway
SImulator. \textit{Bioinformatics}, 22:3067--3074. Disponível em
\url{http://copasi.org}.
\end{thebibliography}
